\section{Related Work}
\label{sec:related}

\subsection{Symbolic and Rule-Based Legal Reasoning}
\label{subsec:related-symbolic}
Formalizing statutory rules as executable logic has a long history in AI and law, beginning with Sergot et al.'s encoding of the British Nationality Act as a Prolog program \cite{sergot1986}, which showed that a hand-authored rule structure could reproduce legal determinations exactly once the relevant facts were supplied - the same separation of fact-finding from rule application that underlies the PYTHEN pipeline in this paper. PROLEG \cite{satoh2010proleg} develops this logic-programming tradition further: a Prolog-based reasoning system that formalizes the Japanese "presupposed ultimate fact" theory (\textit{Yoken-jijitsu-ron}) governing civil litigation, it represents rules as a positive-condition part together with explicit \textit{exception} predicates, so that it can still reach a defensible conclusion under Japanese civil code's burden-of-proof structure when some facts remain unknown, rather than requiring every condition to be resolved to true or false. This explicit treatment of exceptions and open facts goes beyond what a single, fixed Boolean rule tree such as PYTHEN's can express, and marks one concrete direction - alongside deeper rule hierarchies and sequential gating - in which the rule-aware error attribution methodology proposed in this paper could be extended to richer rule structures (Section~\ref{subsec:limitations}). Subsequent work extended this tradition to formal argumentation and defeasible reasoning more broadly, providing frameworks for handling conflicting rules, exceptions, and burdens of proof that a purely Boolean rule tree cannot express \cite{benchcapon2007}. Prakken and Sartor's survey of this line of work from an argumentation perspective \cite{prakken2020} notes that its central appeal, and also limitation, is that soundness and explainability depend on a correctly and completely specified logical formalization, an assumption that traditionally required both the rules and the facts to be entered by hand rather than extracted from unstructured natural-language text.

\subsection{Empirical Legal NLP and Large Language Models}
\label{subsec:related-llm}
The rise of LLMs has shifted attention toward automatically bridging this gap between free text and formal input. LegalBench \cite{guha2023legalbench} introduced a benchmark of 162 tasks - including the \texttt{hearsay} classification task used in this paper - constructed to measure legal reasoning capabilities directly relevant to practicing lawyers, and evaluated LLMs primarily via direct or chain-of-thought prompting for a final label. This evaluation paradigm inherits a known reliability problem: an analysis of chain-of-thought faithfulness \cite{turpin2023unfaithful} shows that such explanations can be systematically unfaithful to the process that actually produced a model's answer, so a fluent justification does not guarantee that the stated reasoning caused the stated conclusion. A related faithfulness study of LLMs used as solvers and autoformalizers in legal reasoning \cite{wang2026knowyourlimits} documents an analogous divergence, finding that benchmark accuracy and logical faithfulness diverge, and identifying failure patterns - such as reporting conclusions unsupported by the model's own formal derivation - that aggregate accuracy does not surface. This paper's central concern, that predicate-level F1 does not indicate decision-level reliability, is a structural analogue of this same divergence, applied to a neuro-symbolic pipeline rather than to free-form or autoformalized reasoning.

\subsection{Hybrid and Neuro-Symbolic Approaches}
\label{subsec:related-hybrid}
Closer to this paper's architecture, a growing body of work couples LLMs with an explicit symbolic or structured component rather than relying on the LLM's final-answer judgment alone. Text-guided legal knowledge graph reasoning \cite{li2021legalkg} casts legal provision prediction as reasoning over a knowledge graph guided by text representations, combining neural text understanding with structured graph inference. A multi-agent framework for verifiable legal reasoning \cite{sadowski2025verifiable} similarly separates knowledge extraction - formalizing statutes into verifiable intermediate representations - from a second phase that applies this formalized knowledge through symbolic inference, reporting substantial accuracy gains over direct prompting on tax-calculation tasks. This body of work establishes that routing an LLM's output through a symbolic or structured intermediate representation is broadly beneficial, consistent with what the neuro-symbolic pipeline in this paper also finds. However, these evaluations, like the aggregate LegalBench baselines above, measure only whether the final decision is correct, not which of the LLM's intermediate extraction errors were responsible when it is not. None of this prior work asks whether a given extraction error would have changed the final decision had it been corrected in isolation, which is the question this paper addresses through predicate-level intervention on a fixed, auditable rule tree.
